\documentclass{article}
\usepackage{graphicx}
\usepackage{amsmath}

\title{Quantum Tomography for High Energy Physics}
\author{N. Viaux, A. Norambuena, F. Aguilar-Carvajal}
\date{May 2026}

\begin{document}

\maketitle

\section{Quantum Tomography at work}
	\subsection{One Spin-1/2 system}
		Let's describe the one spin-1/2 density matrix:
		\begin{equation}
			\rho = \frac{1}{2} \left( \mathcal{I} + \vec{B} \cdot \vec{\sigma} \right)
		\end{equation}

		If we want to analyze a decay $\tau / \overline{\tau} \rightarrow l^+ / l^-$, then the option of meassure just two directions (up/down) isn't available, instead we need to meassure in a space, in a solid angle space. So, we have the intention of meassure in a specific direction in which the product of the decay goes on, say $\hat{n}$, then the proyector operator will take the form:
		\begin{equation}
			\Pi_{\hat{n}} = \frac{1}{2} \left( \mathcal{I} + \hat{n} \cdot \vec{\sigma} \right)
		\end{equation}

		The probability that the decay product be emitted in the $\hat{n}$ direction is:
		\begin{equation}
			\begin{split}
				\text{p}(\hat{n}) & = \text{Tr} \left[ \frac{1}{2} \left( \mathcal{I} + \vec{B} \cdot \vec{\sigma} \right) \frac{1}{2} \left( \mathcal{I} + \hat{n} \cdot \vec{\sigma} \right) \right] \\
				& = \text{Tr} \left[ \frac{1}{4} \left( \mathcal{I} + \vec{B} \cdot \vec{\sigma} + \hat{n} \cdot \vec{\sigma} + \underbrace{(\vec{B} \cdot \vec{\sigma})(\hat{n} \cdot \vec{\sigma})}_{(\vec{B} \cdot \hat{n})I + i(\vec{B} \times \hat{n})\cdot \vec{\sigma}} \right) \right] \\
				& = \frac{1}{4} \text{Tr} \left[ \mathcal{I} + \vec{B} \cdot \vec{\sigma} + \hat{n} \cdot \vec{\sigma} + (\vec{B} \cdot \hat{n})I + i(\vec{B} \times \hat{n})\cdot \vec{\sigma} \right] \\
				& = \frac{1}{4} \text{Tr} \left[ \mathcal{I} + \left( \vec{B} + \hat{n} + i(\vec{B} \times \hat{n}) \right) \cdot \vec{\sigma} + (\vec{B} \cdot \hat{n}) \mathcal{I} \right] \\
				& = \frac{1}{4} \left[ \underbrace{\text{Tr} \left[ \mathcal{I} \right]}_{4} + \underbrace{\text{Tr} \left[ \left( \vec{B} + \hat{n} + i(\vec{B} \times \hat{n}) \right) \cdot \vec{\sigma} \right]}_{\vec{\sigma} \text{'s are traceless}} + \text{Tr} \left[ (\vec{B} \cdot \hat{n})I \right]  \right] \\
				& = \frac{1}{4} \left[ 2 + 2\left( \vec{B} \cdot \hat{n} \right) \right]
			\end{split}
		\end{equation}

		So, the probability is given by:
		\begin{equation}
			\text{p}(\hat{n}) = \frac{1}{2} \left[ 1 + \vec{B}\cdot \hat{n} \right]
		\end{equation}

		However, the probability isn't normalized in the solid angle space. Let's say that $\text{P} (\hat{n}) \propto 1 +\vec{B}\cdot \hat{n}$, then:
		\begin{align}
			\int d\Omega \ \text{p} (\hat{n}) = 1 \\
			\int d\Omega \left[ c (1+\vec{B}\cdot\hat{n}) \right] = 1 \\
			c \underbrace{\int d\Omega}_{4\pi} + c\vec{B} \cdot \underbrace{\int d\Omega \ \hat{n}}_{0} = 1 \\
			c = \frac{1}{4\pi}
		\end{align}

		Finally, the probability density is given by:
		\begin{equation}
			\text{p} (\hat{n}) = \frac{1}{4\pi} \left[ 1 + \vec{B} \cdot \hat{n} \right]
		\end{equation}


\end{document}
